\documentclass[11pt]{article}
\usepackage[a4paper,margin=24mm]{geometry}
\usepackage{amsmath,amssymb,amsthm,booktabs,longtable}
\newtheorem{theorem}{Theorem}\newtheorem{lemma}[theorem]{Lemma}
\newcommand{\PSL}{\operatorname{PSL}}\newcommand{\SL}{\operatorname{SL}}\newcommand{\GL}{\operatorname{GL}}
\newcommand{\PG}{\operatorname{PG}}\newcommand{\E}{\mathrm E}
\title{Parabolics, Borels and residue kernels\ in the $\PSL_3(4)$ catalogue}
\author{L100 literal structural certificate}\date{}
\begin{document}\maketitle
This note identifies the actual subgroups printed in the 29 diagrams of
Gottschalk and Leemans' fixed Halle preprint \cite{GL}, pp.~7--13.
The acquired complete coset catalogue and its source-node bijections are
inputs. We do not repeat the subgroup or geometry census. The retained
21-point realization of $G=\PSL_3(4)$, its order 20160, all its literal
permutations, and its verified coordinate identification with $\PG(2,4)$
are explicit premises. The journal edition is not collated here.

For each catalogue geometry, $H_i$ is the saved stabilizer of its base
object of type $i$, $B=\bigcap_iH_i$, and $K_i$ is the already certified
kernel of $H_i$ on that object's entire residue. These are literal subsets
of the same ordered list of elements of $G$. There are 99 parabolic
positions, 99 kernel positions and 29 Borels. Identical subsets, rather
than equal orders or isomorphic groups, reduce these 227 positions to
77 atoms. Every source position retains its exact atom identifier.
The certificate first checks that a stored generating list generates
exactly its asserted subset.

\begin{lemma}\label{projective}
The actual point or line stabilizers have structure $2^4:A_5$.
The actual unordered point-pair or line-pair stabilizers used here have
structure $2^4:S_3$, and an incident point-line stabilizer has structure
$2^4:A_4$. Each colon is paid by an actual complement.
\end{lemma}
\begin{proof}
Choose coordinates with the fixed point represented by $e_1$.
An element of its stabilizer has an $\SL_3(4)$ representative
$\left(\begin{smallmatrix}a&u\\0&C\end{smallmatrix}\right)$ with
$a\det C=1$. Multiplication by the scalar $a^{-1}$ is permitted because
$a^3=1$. It gives a unique representative with upper-left entry 1 and
lower block in $\SL_2(4)$. Consequently this group is the split extension
of the additive row space $\mathbb F_4^2$ by $\SL_2(4)$.
The latter acts faithfully on its five projective points and is $A_5$.
The dual calculation gives the line stabilizer. A flag imposes one fixed
point in this five-object action, yielding $A_4$. A pair imposes an
unordered two-set, whose stabilizer in $A_5$ is $S_3$.

These descriptions are joined to each actual $H_i$ without assuming that
coordinate and permutation labels agree. The check finds its invariant
point, line, flag, or unordered pair in the retained projective model.
For a fixed line $\ell$, it defines $N$ as the elements fixing all five
points of $\ell$, and $L$ as the stabilizer within $H_i$ of a specified
point outside $\ell$. Dually it uses a pencil and a line not through its
centre. It verifies $|N|=16$, a commuting involution basis of length four,
$N\lhd H_i$, $N\cap L=1$, and every element has a unique expression
$nl$. The faithful five-object action of $L$ is checked on every element;
it is the entire $A_5$, its specified point stabilizer, or its specified
two-set stabilizer. This proves the claimed splitting and the actual
action, in addition to the abstract group name.
\end{proof}

\begin{lemma}\label{geometric}
The other large atoms are $A_6$, $A_5$ and $L_3(2)$, with the indicated
faithful geometric actions.
\end{lemma}
\begin{proof}
For each order-360 atom the checker finds a retained invariant hyperoval
and checks all 360 induced permutations of its six points. They are
distinct and even, hence are precisely $A_6$.
For each order-60 atom it checks either an invariant oval or an actual
five-line orbit. The 60 induced permutations are distinct and even,
proving $A_5$ also for the dual-oval stabilizers. A point action is not
silently substituted for a line action.

For each order-168 atom it finds a retained Baer subplane, checks the
seven points and its seven triple lines, and verifies the unique-line
axiom on every point pair. Every recorded action preserves all seven
lines and is faithful. An automorphism of the Fano plane is determined
by an ordered noncollinear triple: the other four points are recovered
by line intersections. Thus its order is at most $7\cdot6\cdot4=168$.
The exhibited 168 permutations are its full automorphism group
$\GL_3(2)=\PSL_3(2)$, as asserted.
\end{proof}

\paragraph{The small structures.}
All remaining atoms have order at most 72. Their certificates use the
following complete normal forms; no small-group identifier or order-only
recognition is used. Here $D_8$ and $Q_8$ both have order eight.
An elementary abelian atom has a commuting prime-order basis and every
coordinate word exactly once; a cyclic atom has an element of full order.
For $A_4$ and $S_4$, conjugation on the four actual cyclic subgroups of
order three is faithful and gives respectively all 12 even or all 24
permutations. The groups $S_3$ and $7:3$ have actual generators $a,b$,
orders $(3,2)$ or $(7,3)$, and $a^b=a^2$ or $a^4$ as appropriate;
all words $a^ib^j$ occur exactly once.

For $D_8$ use $|a|=4$, $|b|=2$, $a^b=a^{-1}$.
For $Q_8$ use $|a|=|b|=4$, $a^2=b^2$ and $a^b=a^{-1}$.
In both cases the eight words $a^ib^j$ are checked.
The atoms $2\times D_8$ and $2\times Q_8$ additionally exhibit an
involution $z$ outside this eight-element subgroup which commutes with
both generators; all sixteen words are distinct and exhaust the atom.

The three groups $3:S_3$, $3^2:4$ and $3^2:Q_8$ have an explicitly
normal translation group $N=\langle a,b\rangle\cong\mathbb F_3^2$.
The saved complement has order 2, 4 or 8, intersects $N$ trivially,
and its products with the nine coordinate words exhaust the group.
Its conjugation matrices on $(a,b)$ are checked literally. For order 2
the action is $-I$. For order 4 its generator squares to $-I$; its
minimal polynomial $X^2+1$ is irreducible over $\mathbb F_3$, giving
the specified cyclic-four action. For order 8 the matrices are exactly
\[
 \{I,-I\}\ \cup\
 \{M\in\SL_2(3):\operatorname{tr}(M)=0\},                 \tag{1}
\]
the standard quaternion subgroup. This identifies the semidirect
product action as well as its factors. The inversion extension of
$3^2$ can also be written $3:S_3$: choose one cyclic translation factor
as the normal subgroup and the other together with the inverter as
the $S_3$ complement.

\begin{theorem}\label{names}
All 29 Borel names and all 99 parabolic names in the fixed preprint agree
with the actual catalogue groups. All boxed kernel structures agree as
well. The outer colon in diagram 2.3 requires the nonsplit qualification
in Lemma~\ref{kernel}; no group or geometry count changes.
\end{theorem}
\begin{proof}
The preceding witnesses certify all 77 literal atoms. Independently
transcribed names at each printed node are compared through the accepted
source-node bijection, including nodes of equal size. The Borel names
are listed below; the full 99-node table, with each actual parabolic,
kernel, atom and extension certificate, is exported as
\texttt{source\_names.json}. It is a literal finite join, not a new
classification by subgroup orders.
\begin{center}\small
\begin{tabular}{@{}ll@{\qquad}ll@{\qquad}ll@{}}
\toprule Diagram&$B$&Diagram&$B$&Diagram&$B$\\\midrule
2.1&$2^4:A_4$&3.1&$A_4$&4.1,4.8&$3$\\
2.2,2.5&$S_4$&3.2&$4$&4.2--4.7&$2$\\
2.3&$3^2:4$&3.3,3.4&$D_8$&4.9,4.10&$2$\\
2.4&$7:3$&3.5&$3$&4.11&$1$\\
2.6,2.7&$Q_8$&3.6&$S_3$&5.1--5.4&$1$\\
&&3.7&$A_4$&&\\\bottomrule
\end{tabular}
\end{center}
\end{proof}

\begin{lemma}\label{kernel}
Every recorded residue kernel is checked inside its actual parabolic.
The split, direct and nonsplit extensions have the following explicit
interpretations.
\end{lemma}
\begin{proof}
The acquired kernel is rechecked to be contained and normal. For each
nontrivial kernel the checker constructs its entire coset partition
inside $H_i$ and the literal quotient multiplication table. Identical
pairs $(H_i,K_i)$ share one certificate; 61 such pairs cover all 99
positions. A trivial kernel uses the identity map, without manufacturing
a new quotient computation.

In diagram 2.1 the boxed $2^4$ is exactly the translation radical of
Lemma~\ref{projective}. In diagrams 2.2, 2.5, 3.1, 3.3, 3.4 and 3.7
the boxed $2^2$ is a specific subspace of that radical. Its exact basis,
a vector-space complement, and the actual Levi action are retained.
Thus the unparenthesized printed expression is paid in the form
$(2^2\times2^2):L$, with the action on the \emph{whole} radical.
We do not assert a central direct factor of $H_i$ or an invariant
complement merely from this typography. The checker explicitly records
whether its chosen vector complement is invariant.

In diagram 2.7 the actual $Q_8$ kernel has a commuting involution
complement, proving the displayed direct product. In 3.2 the actual
cyclic-four kernel has an elementary-four complement, proving $4:2^2$.
In 3.5 the actual cyclic-three kernel has a disjoint $S_3$ complement,
proving $3:S_3$. These conclusions use the specified kernels, not just
some isomorphic normal subgroups.

In the central nodes of 4.2--4.7 the boxed involution is central and
every square belongs to it. The quotient is elementary abelian of
order eight. The extension is nonsplit: a complement would commute
with the central kernel and make the whole group elementary abelian,
whereas its exhibited dihedral factor has elements of order four.
This is consistent with the source's dotted $2.2^3$ notation.

Finally diagram 2.3 has $H=3^2:Q_8$ and the actual kernel
$K=3^2:C_4$ of index two. Every element outside $K$ is checked to
have order four; in particular none is an involution. Hence there
is no complement of order two. This also follows symbolically by
projecting to $Q_8$: an involution maps to $1$ or its unique central
involution, both in the cyclic-four image of $K$.
The preprint's displayed outer colon in $(3^2:4):2$ must therefore
be read with a nonsplit extension dot in place of that split colon.
The corrected assertion is $(3^2:4).2=3^2:Q_8$ with the specified
kernel. It is an auxiliary structural-notation correction, not a
refutation of the principal geometry classification.
\end{proof}

\paragraph{Verification and scope.}
All products above are products of the actual retained permutations.
A separate witness validator checks the saved maps, normal forms,
geometric factors, quotient tables and actions without searching for
new witnesses. Damaged generators, action tables, kernels, complements
and source associations are rejected. The development record retains
the initial syntax failure and the unsuccessful point-only search for
the dual-oval atom; the final edition handles the five-line action.
The complete ambient list, the coset models, the kernel meaning and
the source-node transport remain acquired inputs, with their earlier
proofs and distinct reviews. This note supplies their missing internal
structural identifications, not another ambient census or a new
measurement. Root's earlier suggestions concerning geometric
stabilizers were advisory; this component's literal code and proof are
authored by the publication-recovery agent and await distinct review.

\begin{thebibliography}{9}
\bibitem{GL} H.~Gottschalk and D.~Leemans,
\emph{Geometries for the group $\PSL(3,4)$},
Halle preprint 99--30, fixed 22-page source, pp.~4, 7--13.
\bibitem{Atlas} J.~H.~Conway et al., \emph{Atlas of Finite Groups},
Oxford University Press, 1985. Notation antecedent cited by \cite{GL};
no Atlas subgroup classification is used to identify these atoms.
\end{thebibliography}
\end{document}
